\section*{Chapter \ref{ch:iv1:timed-test-2} --- Timed Test 2}
\markboth{\protect\omTocMark{}{Answers: Timed Test 2}}{\protect\omTocMark{}{Answers: Timed Test 2}}

\begin{solution}{iq:iv1:timed-test-2:1}
The average is $\frac74$ balls.

Use indicators\index{indicator variables}. A red ball is left over only if all 3 blue balls come out before it. Among that red ball and the 3 blue ones, every order is equally likely, so the chance that the red comes last is $\frac14$. Similarly, a blue ball is left over only if all 5 reds come before it, with chance $\frac16$.

Adding up, $5 \times \frac14 + 3 \times \frac16 = \frac54 + \frac12 = \frac74$.

\iqwrong{Conditioning on which colour runs out and then trying to find the number remaining in each case. It can be done, but it takes much longer, and the case weights are easy to get wrong.}
\end{solution}

\begin{solution}{iq:iv1:timed-test-2:2}
The average is $19.5$ rolls.

A roll gives 7 with chance $\frac16$ and 11 with chance $\frac{1}{18}$. Let $A$ and $B$ be the waiting times for each. The time to see both is $\max(A, B) = A + B - \min(A, B)$\index{inclusion--exclusion}.

$\mathbb{E}A = 6$ and $\mathbb{E}B = 18$. The minimum is the wait for either total, which has chance $\frac16 + \frac{1}{18} = \frac29$ per roll, so $\mathbb{E}\min = 4.5$. The answer is $6 + 18 - 4.5 = 19.5$.

Note that the result is only $1.5$ above the rarer event alone, as the common event usually happens on the way.

\iqwrong{Adding 6 and 18. The two waits overlap, as the rolls count for both at once.}
\end{solution}

\begin{solution}{iq:iv1:timed-test-2:3}
The price is $\frac{20}{3} \approx 6.67$, and the hedge is $\frac23$ of a share.

The call pays 20 or 0. A position of $\Delta$ shares moves by $30\Delta$ between the two states, so to match the 20 difference, $\Delta = \frac23$\index{replication}.

With $\frac23$ of a share, the portfolio is worth 80 or 60. Borrowing 60 at no interest gives exactly 20 or 0, so the call costs $\frac23 \times 100 - 60 = 6.67$.

A second route: the risk-neutral chance of the up state solves $120q + 90(1 - q) = 100$, so $q = \frac13$, and the price is $\frac13 \times 20$. The real chance of the up state appears nowhere.

\iqfollow{What if the call were quoted at 7?}{Sell it, buy $\frac23$ of a share, and borrow the remainder. The position then earns the 33-cent difference in both states, with no risk.}
\end{solution}

\enlargethispage{-3\baselineskip}%
\begin{solution}{iq:iv1:timed-test-2:4}
The probability is $\frac16$.

The boundary is $\sqrt{y} = 1 - \sqrt{x}$, that is $y = (1 - \sqrt{x})^2$, so the event is the region under this curve in the unit square. Its area is
\[
\int_0^1(1 - \sqrt{x})^2\,dx = \int_0^1(1 - 2\sqrt{x} + x)\,dx = 1 - \frac43 + \frac12 = \frac16.
\]
A second route\index{change of variables}: put $U = \sqrt{X}$ and $V = \sqrt{Y}$. Since $P(U \le u) = u^2$, each has density $2u$ on $[0, 1]$, and the answer is $\iint_{u + v < 1}4uv\,du\,dv = 4 \cdot \frac{1}{24} = \frac16$.

\iqwrong{Answering $\frac12$, by treating the boundary as the straight line $x + y = 1$. On $[0, 1]$ we have $\sqrt{x} \ge x$, so the event forces $X + Y < 1$ and covers only part of that triangle: the curve bends in towards the origin.}
\end{solution}

\begin{solution}{iq:iv1:timed-test-2:5}
The probability is $\binom{10}{5}^2/\binom{20}{10} = 63\,504/184\,756 \approx 0.344$.

Every particular sequence of 20 fair flips has the same chance, so given 10 heads in all, each of the $\binom{20}{10}$ placements of those heads is equally likely. Exactly 5 in the first 10 flips means choosing 5 places there and 5 places in the last 10, which gives $\binom{10}{5}^2$ placements\index{hypergeometric distribution}.

For a value, $\binom{2n}{n} \approx 4^n/\sqrt{\pi n}$ turns the ratio into $\sqrt{10\pi}/(5\pi) \approx 0.357$, which is close to the exact $0.344$.

\iqwrong{Answering $\binom{10}{5}/2^{10} \approx 0.246$, the chance of 5 heads in 10 flips with no condition. Fixing the total ties the two halves together: more heads early forces fewer late, which narrows the count in the first half.}

\iqfollow{How do the two compare in general, for large samples?}{The conditional chance is larger by a factor of about $\sqrt{2}$, since the variance of the count in the first half is halved. Here $0.344/0.246 = 1.40$.}
\end{solution}

\begin{solution}{iq:iv1:timed-test-2:6}
The expected number of flips is $\frac83$.

Let $N$ be the number of flips, so $P(N = k) = 2^{-k}$ and $P(N = 2j) = 4^{-j}$. These sum to $\frac13$, the chance of an even count. Given an even count, $P(N = 2j) = 3 \cdot 4^{-j}$, so $N/2$ is geometric with success chance $\frac34$ and mean $\frac43$\index{conditioning on an event}. So $N$ has conditional mean $\frac83$.

A direct sum agrees: $\sum_j 2j\,4^{-j} = \frac89$, and $\frac89 \big/ \frac13 = \frac83$.

\iqwrong{Saying 4, by doubling the unconditional mean of 2, since $N$ is twice $N/2$. Given an even count, $N/2$ is not distributed like $N$: the value 1 has chance $\frac34$, not $\frac12$.}

\iqfollow{What does this give for a coin with heads chance $p$?}{$2/(1 - (1 - p)^2)$. It tends to 2 as $p$ tends to 1, as it must, since 2 is the smallest even count.}
\end{solution}

\begin{solution}{iq:iv1:timed-test-2:7}
The probability is $\frac{12}{38} = \frac{6}{19} \approx 0.316$.

The 39 cards you cannot see form three hands of 13. Place the two missing aces one at a time among those 39 slots\index{sequential placement}. The first goes anywhere. The second must take one of the 12 other slots in that hand, out of the 38 left.

A count of pairs gives the same value: of the $\binom{39}{2} = 741$ pairs of slots, $3\binom{13}{2} = 234$ lie in one hand, and $\frac{234}{741} = \frac{6}{19}$.

\iqwrong{Saying $\frac13$, since there are three hands for the second ace. A hand already holding an ace has only 12 free slots, against 13 in each of the others. With hands of $h$ cards the chance is $(h - 1)/(3h - 1)$: this is 0 for $h = 1$, and it tends to $\frac13$ only for large $h$.}
\end{solution}

\begin{solution}{iq:iv1:timed-test-2:8}
The average is $\frac{20}{19} \approx 1.05$ couples.

Use an indicator\index{indicator variables} for each couple. Fix one partner's seat. The other partner is equally likely to be in any of the remaining 19 seats, and exactly 2 of them are adjacent. So each couple sits together with chance $\frac{2}{19}$.

By linearity, the mean is $10 \times \frac{2}{19} = \frac{20}{19}$. The indicators are dependent, but linearity does not care.

The mean is $2n/(2n - 1)$ for $n \ge 2$ couples, which tends to 1 as $n$ grows, and the count tends to a Poisson with mean 1.

\iqwrong{Trying to find the distribution of the number of adjacent couples. That needs inclusion--exclusion, while the mean needs only one line.}
\end{solution}

\begin{solution}{iq:iv1:timed-test-2:9}
$\rho = 0.24/0.99 \approx 0.242$.

The variance of an equally weighted basket of $N$ names is\index{implied correlation}
\[
\sigma_I^2 = \sigma^2\left(\frac1N + \left(1 - \frac1N\right)\rho\right).
\]
With $\sigma_I^2/\sigma^2 = (15/30)^2 = 0.25$ and $N = 100$: $0.25 = 0.01 + 0.99\rho$, so $\rho = 0.24/0.99$.

Without the correction for finite $N$, the quick answer is $0.25$, which is close. For large $N$, the ratio of the two variances is just the correlation, and that is how a desk reads an average correlation from an index and its members.

Sanity check: at $\rho = 1$, the index vol is 30\%, and at $\rho = 0$, it is 3\%.

\iqfollow{The index's options imply 15\%, and the stocks' options imply 30\%. A trader thinks that realised correlation will be $0.15$. What vol does this imply for the index?}{$30\% \times \sqrt{0.01 + 0.99 \times 0.15} = 11.9\%$, below the 15\% implied. That is the trade that dispersion desks consider.}
\end{solution}

\begin{solution}{iq:iv1:timed-test-2:10}
The probability is $\frac{5 - \pi}{4} \approx 0.4646$.

The ratio rounds to 0 when $X/Y < \frac12$, with probability $\frac14$ (a triangle of base $\frac12$ and height 1).

For $r \ge 1$, $P(X/Y > r) = 1/(2r)$. So the ratio rounds to $2k$, for $k \ge 1$, with probability
\[
\frac12\left(\frac{1}{2k - \frac12} - \frac{1}{2k + \frac12}\right) = \frac{1}{4k - 1} - \frac{1}{4k + 1}.
\]
The sum over $k$ is $\frac13 - \frac15 + \frac17 - \frac19 + \dots$, which by the Leibniz series\index{Leibniz series} is $1 - \frac{\pi}{4}$.

Adding $\frac14$ gives $\frac54 - \frac{\pi}{4}$. A simulation returns $0.465$. It is below $\frac12$, because the odd value 1 alone takes about $0.42$ of the mass, more than the even value 0.

\iqwrong{Using $P(X/Y > r) = 1/r$. The region $X > rY$ in the unit square is a triangle of base 1 and height $1/r$, so there is a factor of $\frac12$.}
\end{solution}

\begin{solution}{iq:iv1:timed-test-2:11}
The probability is about $0.736$.

Split by the score at which the game ends. With $p = 0.6$ and $q = 0.4$:

Winning to love, to 15 and to 30: $p^4(1 + 4q + 10q^2) = 0.1296 \times 4.2 = 0.544$.

Reaching deuce, 3--3: $\binom63p^3q^3 = 20 \times 0.216 \times 0.064 = 0.276$.

From deuce, two points decide or return to deuce. Win two with $p^2$, lose two with $q^2$, so the chance of winning from deuce is\index{geometric series}
\[
\frac{p^2}{p^2 + q^2} = \frac{0.36}{0.52} = 0.692.
\]
Total: $0.544 + 0.276 \times 0.692 = 0.736$. An exact recursion over all scores agrees.

A $0.6$ point edge becomes a $0.74$ game edge, and that is how small edges compound through repeated play.

\iqwrong{Forgetting the coefficients, or writing $\binom{6}{3}$ wrongly for love, 15 and 30. The last point must be won, so the counts are $\binom{3}{0}$, $\binom{4}{1}$ and $\binom{5}{2}$.}
\end{solution}

\begin{solution}{iq:iv1:timed-test-2:12}
The sum is 6.

Let $N$ be the number of flips until the first head, so $P(N = n) = 2^{-n}$. The sum is $\mathbb{E}N^2$\index{geometric distribution}. For a geometric variable with $p = \frac12$, the mean is 2, and the variance is $(1 - p)/p^2 = 2$. So $\mathbb{E}N^2 = 2 + 4 = 6$.

A direct route: from $\sum x^n = 1/(1 - x)$, differentiating and multiplying by $x$ twice gives $\sum n^2x^n = x(1 + x)/(1 - x)^3$. At $x = \frac12$, this is $\frac12 \cdot \frac32 \cdot 8 = 6$. A partial sum confirms it.

\iqfollow{What is $\sum n^3/2^n$?}{26. Differentiating once more gives $x(1 + 4x + x^2)/(1 - x)^4 = \frac12 \cdot \frac{13}{4} \cdot 16 = 26$.}

\iqwhy{Turning a series into a moment saves algebra and it gives a second way of checking the answer.}
\end{solution}
